% Based on templates/tutorial.tex; finished notes omit conversational checks.
% Source: tut10.pdf, p. 10-1 (PDF p. 1), Questions 1--3. No answer key supplied.
% Q1 interprets each eviction entry as one page reference.
% Q2b uses a four-reference window, including the current reference, and keeps
% only its distinct pages resident. Initial window assumes consecutive times 160--163.
% Session: 2026-10-27, teaching Week 11. Replacement traces independently checked.

\begingroup
\setlength{\parskip}{0.4em}
\renewcommand{\recordid}[1]{}
\renewcommand{\relatedtopic}[2]{\hyperref[topic:#1]{#2}}

\clearpage
\exercisegroup{SC2005-TUT10}{Tutorial 10：LRU 反推、工作集与抖动}
\label{tutorial:SC2005-TUT10}

{\small 2026-10-27（第十一教学周）。来源：\texttt{tut10.pdf}，\emph{Virtual Memory (Part 2)}，p.~10-1（PDF p.~1），Questions 1--3。原文件无参考答案；以下为按题意推导并逐步核对的解答。}

\exercisequestion{SC2005-TUT10-Q01}{根据换出记录反推页面引用串}

\textbf{对应讲义：}\relatedtopic{SC2005-L18-T03}{LRU 页面置换}。

\textbf{题目。}请求分页系统给进程分配 $4$ 个页框，初始全部装有页面；进程只访问页面 $1,2,3,4,5$，使用 LRU。给定依次换出的页面为
\[
  3,\;1,\;2,\;4,\;5,\;N,\;1,
\]
其中 $N$ 表示没有页面被换出。求 (a) 初始装入的页面；(b) 对应的页面引用串。按题意将换出记录的每一项对应到一次页面引用。

\begin{checkedsolution}
\textbf{(a) 初始页面为 $\{1,2,3,4\}$，LRU 顺序为 $3\to1\to2\to4$。}顺序从最久未访问排到最近访问，并非物理页框编号。前四次都是缺页；新装入的页面刚被访问，原有且尚未访问的页面会先被淘汰，因此前四个换出页就是初始的四个页面及其 LRU 顺序。

\textbf{(b) 每次缺页访问的，都是五个页面中唯一不在内存的那个。}最初只缺页 $5$，所以先访问 $5$；换出 $3$ 后只缺 $3$，下一次就访问 $3$。依此得到前五次引用 $5,3,1,2,4$。

此时 LRU 顺序为 $[3,1,2,4]$。第六次没有换出页；由于页框一直满载，这次必为命中。要使第七次换出 $1$，第六次必须访问 $3$，将它移到最近使用端；随后访问缺失的 $5$。完整过程如下：

\begin{center}
\begin{tabular}{@{}rccc@{}}
\toprule
步 & 引用页 & 换出页 & 访问后的 LRU 顺序（久 $\to$ 近） \\
\midrule
初始 & --- & --- & $[3,1,2,4]$ \\
1 & 5 & 3 & $[1,2,4,5]$ \\
2 & 3 & 1 & $[2,4,5,3]$ \\
3 & 1 & 2 & $[4,5,3,1]$ \\
4 & 2 & 4 & $[5,3,1,2]$ \\
5 & 4 & 5 & $[3,1,2,4]$ \\
6 & 3 & $N$ & $[1,2,4,3]$ \\
7 & 5 & 1 & $[2,4,3,5]$ \\
\bottomrule
\end{tabular}
\end{center}
因此引用串为 $\boxed{5,3,1,2,4,3,5}$。换出串与引用串含义不同；LRU 命中也会更新最近使用顺序。
\end{checkedsolution}

\begin{samepage}
\exercisequestion{SC2005-TUT10-Q02}{固定四页框 LRU 与四次引用窗口的工作集}

\textbf{对应讲义：}\relatedtopic{SC2005-L18-T03}{LRU}；\relatedtopic{SC2005-L19-T03}{Working-set model（工作集模型）}。

\textbf{题目。}进程的四个页框在下一次缺页前具有如下状态：\par
\end{samepage}

\begin{center}
\begin{tabular}{@{}rrrr@{}}
\toprule
页框号 & 页面号 & 装入时间 & 最后访问时间 \\
\midrule
0 & 1 & 60 & 161 \\
1 & 0 & 130 & 160 \\
2 & 3 & 26 & 162 \\
3 & 2 & 20 & 163 \\
\bottomrule
\end{tabular}
\end{center}
剩余页面引用串为
\[
  4,\;0,\;0,\;0,\;2,\;4,\;2,\;1,\;0,\;3,\;2.
\]
分别求缺页次数并展示过程：(a) 固定分配四个页框，使用 LRU；(b) 工作集窗口长度为 $4$，内存只保留工作集中的页面，不再采用固定四页框分配。

\begin{checkedsolution}
\textbf{(a) LRU 按最后访问时间排序。}初始顺序为 $[0,1,3,2]$；装入时间不决定 LRU 的淘汰次序。命中将当前页移到最近使用端，缺页则淘汰最久未访问页并装入当前页。

\textbf{(b) 窗口计访问次数，工作集计不同页面。}按时间 $160$--$163$ 对应连续四次页面引用来初始化窗口 $[0,1,3,2]$。每次访问后取最近四次引用（含本次），重复引用也占窗口位置；去重得到工作集。例如 $[2,4,0,0]$ 对应 $\{0,2,4\}$，只需三个页框。

访问前不在驻留集中的页面产生缺页；访问后按本题规定，移除窗口中已不再出现的页面，即使还有空闲页框也不保留。

下表均记录访问后的状态。LRU 一列按久到近排列，窗口按旧到新排列；F 为缺页，H 为命中，WS 为工作集策略。两个策略从同一初始状态分别运行。

{\small
\begin{longtable}{@{}rrlcllc@{}}
\toprule
步 & 引用 & LRU 顺序 & LRU & 最近四次引用 & 工作集 & WS \\
\midrule
\endhead
\bottomrule
\endfoot
初始 & --- & $[0,1,3,2]$ & --- & $[0,1,3,2]$ & $\{0,1,2,3\}$ & --- \\
1 & 4 & $[1,3,2,4]$ & F & $[1,3,2,4]$ & $\{1,2,3,4\}$ & F \\
2 & 0 & $[3,2,4,0]$ & F & $[3,2,4,0]$ & $\{0,2,3,4\}$ & F \\
3 & 0 & $[3,2,4,0]$ & H & $[2,4,0,0]$ & $\{0,2,4\}$ & H \\
4 & 0 & $[3,2,4,0]$ & H & $[4,0,0,0]$ & $\{0,4\}$ & H \\
5 & 2 & $[3,4,0,2]$ & H & $[0,0,0,2]$ & $\{0,2\}$ & F \\
6 & 4 & $[3,0,2,4]$ & H & $[0,0,2,4]$ & $\{0,2,4\}$ & F \\
7 & 2 & $[3,0,4,2]$ & H & $[0,2,4,2]$ & $\{0,2,4\}$ & H \\
8 & 1 & $[0,4,2,1]$ & F & $[2,4,2,1]$ & $\{1,2,4\}$ & F \\
9 & 0 & $[4,2,1,0]$ & H & $[4,2,1,0]$ & $\{0,1,2,4\}$ & F \\
10 & 3 & $[2,1,0,3]$ & F & $[2,1,0,3]$ & $\{0,1,2,3\}$ & F \\
11 & 2 & $[1,0,3,2]$ & H & $[1,0,3,2]$ & $\{0,1,2,3\}$ & H \\
\end{longtable}
}

LRU 在第 $1,2,8,10$ 步缺页，换出页依次为 $0,1,3,4$；工作集策略在第 $1,2,5,6,8,9,10$ 步缺页。因此
\[
  \boxed{\text{LRU：}4\text{ 次缺页}},\qquad
  \boxed{\text{工作集：}7\text{ 次缺页}}.
\]
本题工作集释放了页框，却增加缺页；不能据此判定它一般劣于 LRU。
\end{checkedsolution}

\par\medskip
\begin{minipage}{\textwidth}
\exercisequestion{SC2005-TUT10-Q03}{从资源利用率判断分页瓶颈}

\textbf{对应讲义：}\relatedtopic{SC2005-L19-T02}{Thrashing（抖动）}；\relatedtopic{SC2005-L19-T03}{工作集需求}。

\textbf{题目。}某请求分页系统测得 CPU 利用率为 $20\%$，分页磁盘利用率为 $97.7\%$，其他 I/O 设备利用率为 $5\%$。判断下表各措施是否可能提高 CPU 利用率。\par
\end{minipage}

\begin{checkedsolution}
CPU 较空闲而分页磁盘接近满负荷，强烈提示分页 I/O 瓶颈，符合抖动的典型表现。改善方向是减少缺页或缩短调页等待，而不是仅提高 CPU 算力或磁盘容量。

{\small
\begin{longtable}{@{}p{0.4cm}>{\raggedright\arraybackslash}p{3.45cm}>{\raggedright\arraybackslash}p{1.7cm}>{\raggedright\arraybackslash}p{\dimexpr\textwidth-5.55cm-6\tabcolsep\relax}@{}}
\toprule
项 & 措施 & 判断 & 原因或条件 \\
\midrule
\endhead
\bottomrule
\endfoot
a & 更快的 CPU & 通常无助 & 未减轻分页 I/O 等待，CPU 算力并非当前瓶颈。 \\
b & 容量更大的分页磁盘 & 通常无助 & 增加容量不等于加快读写，也不增加可驻留的内存页面。 \\
c & 增加多道程序程度 & 通常恶化 & 更多进程竞争有限页框，工作集更难同时驻留，可能增加缺页。 \\
d & 降低多道程序程度 & 有助 & 减少同时驻留的进程，让剩余进程获得足够页框。 \\
e & 增加主存 & 有助 & 容纳更多工作集页面，减少缺页及调页等待。 \\
f & 更快的硬盘或多个硬盘 & 有助，有条件 & 更快的分页 I/O 缩短等待；多个磁盘需实际分担分页读写，单纯增加闲置磁盘无效。 \\
g & 增大页面大小 & 取决于局部性 & 良好的空间局部性可能使一次调入覆盖更多后续访问；但同样主存下页框数减少，内部碎片及无用数据传输也可能增加。 \\
\end{longtable}
}

\textbf{主要选择为 d、e、f；g 不能仅凭给出的利用率确定。}更快的磁盘可缓解等待，但不直接解决工作集所需内存不足的问题。
\end{checkedsolution}

\subsubsection*{一分钟回顾}

LRU 的次序由最近访问决定，命中也要更新；反推时先辨清引用页与换出页。工作集窗口长度统计访问次数，驻留页数统计窗口内不同页，两者不必相等。本题 LRU 与工作集缺页数分别为 $4$、$7$。CPU 低利用率同时伴随分页磁盘繁忙时，应优先减少内存竞争、增加主存或加快分页 I/O。

\endgroup
